\documentclass[11pt,a4paper]{article}
\usepackage[margin=25mm]{geometry}
\usepackage[T1]{fontenc}
\usepackage[utf8]{inputenc}
\usepackage{lmodern,microtype}
\usepackage{amsmath,amssymb,booktabs}
\usepackage{hyperref}
\title{Reconstructing the SmartWeld OSLW Physics Model: A Reproducible Forward Model and a Framework for Physics-Informed Machine Learning}
\author{Andrey Ni}
\date{Preprint draft v0.1 --- 30 September 2026}
\begin{document}
\maketitle
\begin{abstract}
This study reconstructs the OSLW (Optimized Single-pass Laser Welding) model
distributed with SmartWeld 3.0 in the 12 March 2012 MATLAB M-files release.
The work addresses reproducibility of legacy scientific software by recovering
the model structure, normalization, coefficients, and numerical behavior. The
recovered forward path includes process-variable normalization, material
parameters, shielding-gas coefficients, penetration depth, energy-transfer
efficiency, melting efficiency, weld width, and penetration sensitivity. An
independent implementation reproduces two historical SmartWeld GUI cases for
304 stainless steel under argon shielding to the printed precision of the
original outputs. The historical inverse workflow is treated separately because
its contour-generation dependency, \texttt{ccontours.m}, has not yet been
recovered into the targeted source artifact. The validated forward model is
proposed as a physics-based prior for design of experiments, surrogate modelling,
constrained optimization, and later experimental research.
\end{abstract}
\section{Introduction}
Legacy scientific software may contain calibrated physical models whose numerical
behavior is difficult to reproduce outside the original executable. Reproducibility
requires explicit mathematical structure, normalization, coefficients, parameter
domain, and selection logic. This study uses SmartWeld 3.0 OSLW as a case study.
\subsection{Objectives}
The work addresses five tasks: recovery of the OSLW forward model; documentation
of normalization and parameter conventions; independent implementation; historical
GUI validation; and definition of a controlled interface for physics-informed ML
and optimization.
\subsection{Research questions}
Can the OSLW forward model be reconstructed unambiguously? Does the independent
implementation reproduce original GUI outputs? Which parts of the inverse workflow
are source-reproducible? Can the recovered model serve as a deterministic physics
prior for ML? What additional evidence is required before extension to DED?
\section{Source Provenance and Software Archaeology}
The investigated release is the 12 March 2012 SmartWeld M-files distribution.
The standalone application was generated with MATLAB 2010a. The recovered OSLW
forward call chain is
\begin{equation}
\boxed{\texttt{app\_specific\_data}\rightarrow
\texttt{schedulermodels}\rightarrow\texttt{queryfunc}}.
\end{equation}
The model depends on global application state, material normalization, gas
selection, spot diameter, and output scaling.
\section{Mathematical Reconstruction}
\subsection{Input and material normalization}
\begin{equation}
s_P=P/1600,\qquad s_V=V/120,\qquad s_d=d/0.0294.
\end{equation}
For 304 stainless steel:
\begin{equation}
\alpha_s=5.18/110,\qquad \Delta h_s=9.41/11.8.
\end{equation}
For argon:
\begin{align}
C_1&=0.9015503591510303,& C_2&=0.6327730407815897,& C_3&=0.2582188708198456,\\
C_8&=29.32835498219138,& C_9&=0.2903225193306445,& C_{10}&=0.3160507149639552.
\end{align}
\subsection{Penetration}
\begin{equation}
p_s=\frac{\alpha_sC_8(s_P/[s_d^{C_{10}}s_V^{C_9}])}{5}.
\end{equation}
\subsection{Energy-transfer efficiency}
\begin{equation}
\eta_{\mathrm{ET},s}=
\frac{C_1-C_2^{\pi/[2\arctan(C_3s_d/p_s)]}}{0.952}.
\end{equation}
\subsection{$Ry$ parameter and melting efficiency}
\begin{equation}
Ry_s=\frac{s_P\eta_{\mathrm{ET},s}s_V}{\alpha_s^2\Delta h_s},
\qquad
Ry=Ry_s\frac{1600\cdot120\cdot0.952}{110^2\cdot11.8},
\end{equation}
\begin{equation}
\eta_m=0.48-0.29e^{-Ry/6.82}-0.17e^{-Ry/58.8}.
\end{equation}
\subsection{Weld width}
The source evaluates
\begin{equation}
A_s=\frac{Ry\eta_m\alpha_s^2}{s_V^2},
\qquad A=A_s(110/120)^2,
\end{equation}
followed by source-defined output scaling in \texttt{queryfunc}.
\subsection{Penetration sensitivity}
\begin{equation}
\frac{\partial p}{\partial P}=
\frac{\alpha_sC_8}{1600s_d^{C_{10}}s_V^{C_9}}.
\end{equation}
\section{Independent Implementation}
The computational sequence is
\begin{equation}
(P,V,d)\rightarrow(s_P,s_V,s_d)\rightarrow
p,\eta_{\mathrm{ET}},Ry,\eta_m,w,\partial p/\partial P.
\end{equation}
No parameter fitting is introduced in this reconstruction.
\section{Validation}
For Case 1, 304 stainless steel, argon, $d=0.0118$ cm:
\begin{equation}
P=685.000000~\mathrm{W},\qquad V=53.082625~\mathrm{mm/s}.
\end{equation}
The historical output reports
\begin{align}
p&=0.999860~\mathrm{mm},&w&=0.625687~\mathrm{mm},\\
\eta_{\mathrm{ET}}&=0.679181,&\eta_m&=0.447787,\\
\partial p/\partial P&=1.459649~\mu\mathrm{m/W}.
\end{align}
For Case 2:
\begin{equation}
P=786.666667~\mathrm{W},\qquad V=52.839095~\mathrm{mm/s},
\end{equation}
with
\begin{align}
p&=1.149791~\mathrm{mm},&w&=0.678073~\mathrm{mm},\\
\eta_{\mathrm{ET}}&=0.718454,&\eta_m&=0.457256,\\
\partial p/\partial P&=1.461599~\mu\mathrm{m/W}.
\end{align}
\begin{table}[ht]
\centering
\caption{Historical SmartWeld OSLW regression cases.}
\begin{tabular}{lcc}
\hline
Quantity & Case 1 & Case 2\\ \hline
Power (W) & 685.000000 & 786.666667\\
Speed (mm/s) & 53.082625 & 52.839095\\
Penetration (mm) & 0.999860 & 1.149791\\
Width (mm) & 0.625687 & 0.678073\\
ETE & 0.679181 & 0.718454\\
Melting efficiency & 0.447787 & 0.457256\\
Sensitivity ($\mu$m/W) & 1.459649 & 1.461599\\ \hline
\end{tabular}
\end{table}
The reconstructed forward model reproduces the historical values to the printed
precision. This is the present numerical validation result.
\section{Historical Inverse Workflow}
The recovered GUI code shows that depth-only mode obtains a process contour and
selects its midpoint before evaluating \texttt{queryfunc}. Thus the historical
GUI solution depends on contour-generation implementation. The exact
\texttt{ccontours.m} source remains the critical unresolved dependency.
\section{Physics-Informed Machine-Learning Framework}
The proposed research cycle is
\begin{equation}
\boxed{\mathrm{DOE}\rightarrow\mathrm{SmartWeld}\rightarrow
\mathrm{dataset}\rightarrow\mathrm{surrogate}\rightarrow
\mathrm{constrained\ optimization}\rightarrow\mathrm{experiment}\rightarrow
\mathrm{update}}.
\end{equation}
The initial benchmark should predict penetration, width, energy-transfer efficiency,
melting efficiency, and penetration sensitivity. Baseline data-driven surrogates
should be compared with physics-informed or physics-constrained alternatives.
Performance claims are reserved for actual benchmark calculations and experiments.
\section{Extension Toward DED}
SmartWeld OSLW is a laser-welding model. The present evidence does not establish
direct SmartWeld variables for powder feed, deposition efficiency, nozzle geometry,
stand-off distance, or layer-to-layer thermal history. Therefore
\begin{equation}
\text{SmartWeld physics prior}+
\text{DED-specific variables}+
\text{experimental data}
\rightarrow\text{validated hybrid model}.
\end{equation}
\section{Limitations}
The exact \texttt{ccontours.m} implementation is not yet recovered; the initial
regression contains two historical GUI cases; the empirical validity domain must
not be exceeded without validation; OSLW is not automatically a DED model; and
ML/autonomous-experiment extensions remain proposed until tested.
\section{Research Program}
Phase I: source completeness and recovery of \texttt{ccontours.m}.\\
Phase II: deterministic numerical reproducibility and inverse reconstruction.\\
Phase III: physics-informed ML, uncertainty estimation and constrained optimization.\\
Phase IV: experimental laser/DED validation.
\section{Expected Contribution}
The intended contribution is a reproducible methodology for recovering a historical
physics-based scientific model and exposing it as a controlled computational prior
for modern data-driven process optimization. The first preprint focuses on source
reconstruction and numerical validation; ML and DED results will be added only when
supported by actual calculations or experiments.
\section{Conclusion}
The SmartWeld OSLW forward model has been reconstructed from historical MATLAB
source and independently reproduced for two historical GUI cases. The inverse
workflow remains partially unresolved because its contour-generation dependency
has not yet been recovered. The forward model provides a reproducible basis for a
subsequent physics-informed ML benchmark, while extension to DED remains a separate
validation problem.
\appendix
\section{Evidence Policy}
Each manuscript statement should be internally classified as Source, Regression,
Reconstruction, Inference, or Proposal. Only Source and Regression claims should
be treated as established facts; Inference and Proposal claims remain qualified.
\end{document}